%! Author = WelJo
%! Date = 4/30/2026

%! Einfügung und Veränderung von "Forschungsstand und theoretische Grundlage" zu Methodik und Vorgehensweise (Bao)

\section{Methodik und Vorgehensweise}
\label{sec:methodik-und-vorgehensweise}

%% Beginning ===========================

„ In diesem Abschnitt ist darzulegen, welche grundlegende Methodik Sie zur
 Beantwortung der Fragestellung einsetzen werden und wie Sie dabei
vorgehen wollen. Beschreiben Sie in diesem Abschnitt die einzelnen
Schritte, die Sie durchführen und wie ein Schritt den nächsten
beeinflusst. Beispielsweise könnten Sie zuerst eine Literaturstudie zu
bestimmten Themen unter einem gewissen Gesichtspunkt durchführen. Darauf
 aufbauend könnten Sie anschließend z. B. ein Konzept ausarbeiten (z. B.
 für eine Studie oder eine softwaretechnische Problemlösung).
Möglicherweise werden Sie Ihr Konzept ganz oder teilweise umsetzen
und/oder theoretisch oder empirisch evaluieren.
Wir wissen, dass Sie
nicht alle Schritte im Exposé bereits detailliert planen können, möchten
 aber, dass Sie bereits vor Beginn der Arbeit die Komplexität Ihres
Themas richtig abschätzen. Mögliche Entscheidungen, die Sie erst während
 der späteren Bearbeitung treffen können, benennen Sie daher bitte auch.
 Bitte beachten Sie, dass die Darstellung der Gliederung der
schriftlichen Arbeit nicht zur Vorgehensweise gehört. Diese wird
nachfolgend separat beschrieben. “

Im ersten Schritt gehört die Einarbeitung in die theoretischen Grundlagen zur Kenntnisgewinnung über den Rate of Accessability nach Kaub et al. und wiederstandsfähigen Infrastrukturen.


%! Von dir der Konzeptteil

%Das Konzept ist das Herzstück deines Exposés: Hier stellst du deine Fragestellung, deine \href{https://www.scribbr.de/methodik/hypothesen-formulieren/}{Hypothese} (Ergebnisannahme) und gegebenenfalls auch deine \href{https://www.scribbr.de/category/methodik/}{Methodik} vor.

%„Im Rahmen der Bachelorarbeit wird eine Analyse dieser \href{https://www.scribbr.de/methodik/fallstudie/}{Case Study} durchgeführt, um zu untersuchen, wie sich der Tweet durch die Medienlandschaft gezogen hat.“




%Ein kleiner (nicht zu ausführlicher) Überblick über die bestehende
%Literatur und – je nach Forschungsgebiet – eine theoretische Grundlage
%für deine Forschung.


%„Wie sich der Two-Step-Flow und das Konzept von Opinion
%Leadership auf Twitter äußern, wurde so noch nicht in der Literatur
%behandelt und der Aspekt der politischen Kommunikation wurde ebenfalls
%weitgehend vernachlässigt.“

%% Ending ==============================